\section{The texts of the Mass}
The Latin follows the Vatican Missal of 1962. The page and marginal number under each heading refer to that edition. For the passages of Scripture, the first English text is the Douay--Rheims Bible in the Challoner revision. The liturgical introductions, responses and minor doxology are from the historical hand missal of Lasance and Walsh. Its rendering of a chant or lesson follows where it differs. The same hand missal supplies the English of the prayers and Preface. The Sunday cue for each prayer and the full ordinary conclusion are distinguished below.

\subsection{Introit}
\label{proper-introit}
\properlocus{No. 1, p. [1], leaf 82, left column. Ps 24(25):1--3, 4.}

\begin{latinproper}
Ad te levávi ánimam meam: Deus meus, in te confído, non erubéscam: neque irrídeant me inimíci mei: étenim univérsi, qui te exspéctant, non confundéntur. Ps. Vias tuas, Dómine, demónstra mihi: et sémitas tuas édoce me. V. Glória Patri, et Fílio, et Spirítui Sancto. Sicut erat in princípio, et nunc, et semper, et in sǽcula sæculórum. Amen.

Ad te levávi ánimam meam: Deus meus, in te confído, non erubéscam: neque irrídeant me inimíci mei: étenim univérsi, qui te exspéctant, non confundéntur.
\end{latinproper}

\begin{englishwitness}{Challoner, with historical liturgical forms}
To thee, O Lord, have I lifted up my soul. In thee, O my God, I put my trust; let me not be ashamed. Neither let my enemies laugh at me: for none of them that wait on thee shall be confounded. Ps. Shew, O Lord, thy ways to me, and teach me thy paths. V. Glory be to the Father, and to the Son, and to the Holy Ghost. As it was in the beginning, is now, and ever shall be, world without end. Amen.

To thee, O Lord, have I lifted up my soul. In thee, O my God, I put my trust; let me not be ashamed. Neither let my enemies laugh at me: for none of them that wait on thee shall be confounded.
\end{englishwitness}

\begin{englishwitness}{Lasance and Walsh, hand-missal English}
To Thee, O Lord, have I lifted up my soul: in Thee, O my God, I put my trust; let me not be ashamed. Neither let my enemies laugh at me; for none of them that wait on Thee shall be confounded. Ps. Show, O Lord, Thy ways to me, and teach me Thy paths. V. Glory be to the Father, and to the Son, and to the Holy Ghost. As it was in the beginning, is now, and ever shall be, world without end. Amen.

To Thee, O Lord, have I lifted up my soul: in Thee, O my God, I put my trust; let me not be ashamed. Neither let my enemies laugh at me; for none of them that wait on Thee shall be confounded.
\end{englishwitness}

\subsection{Collect}
\label{proper-collect}
\properlocus{No. 2, p. [1], leaf 82, right column.}

\begin{latinproper}
Orémus. Excita, quǽsumus, Dómine, poténtiam tuam, et veni: ut ab imminéntibus peccatórum nostrórum perículis, te mereámur protegénte éripi, te liberánte salvári: Qui vivis et regnas cum Deo Patre in unitáte Spíritus Sancti, Deus: per ómnia sǽcula sæculórum. R. Amen.
\end{latinproper}

\begin{englishwitness}{Lasance and Walsh, Sunday text}
Bestir, O Lord, Thy might, we pray Thee and come; That, defended by Thee, we may deserve rescue from approaching dangers brought on by our sins, and being set free by Thee, obtain our salvation. Who livest.
\end{englishwitness}

\begin{englishwitness}{Full conclusion, Prayer Conclusions, p. 819}
Who livest and reignest, with God the Father, in the unity of the Holy Ghost, God, world without end. R. Amen.
\end{englishwitness}

\subsection{Epistle}
\label{proper-epistle}
\properlocus{No. 3, pp. [1]--2, leaves 82--83. Rom 13:11--14a.}

\begin{latinproper}
Léctio Epístolæ beáti Pauli Apóstoli ad Romános. Fratres: Sciéntes, quia hora est iam nos de somno súrgere. Nunc enim própior est nostra salus, quam cum credídimus. Nox præcéssit, dies autem appropinquávit. Abiciámus ergo ópera tenebrárum, et induámur arma lucis. Sicut in die honéste ambulémus: non in comessatiónibus, et ebrietátibus, non in cubílibus, et impudicítiis, non in contentióne, et æmulatióne: sed induímini Dóminum Iesum Christum.
\end{latinproper}

\begin{englishwitness}{Challoner, with historical liturgical forms}
Lesson from the Epistle of blessed Paul the apostle to the Romans. Brethren, knowing the season, that it is now the hour for us to rise from sleep. For now our salvation is nearer than when we believed. The night is passed And the day is at hand. Let us, therefore cast off the works of darkness and put on the armour of light. Let us walk honestly, as in the day: not in rioting and drunkenness, not in chambering and impurities, not in contention and envy. But put ye on the Lord Jesus Christ.
\end{englishwitness}

\begin{englishwitness}{Lasance and Walsh, hand-missal English}
Lesson from the Epistle of blessed Paul the apostle to the Romans. Brethren, knowing the time, that it is now the hour for us to rise from sleep; for now our salvation is nearer than when we believed. The night is past, and the day is at hand; let us therefore cast off the works of darkness, and put on the armor of light. Let us walk honestly as in the day; not in rioting and drunkenness, not in chambering and impurities, not in contention and envy; but put ye on the Lord Jesus Christ.
\end{englishwitness}

\subsection{Gradual}
\label{proper-gradual}
\properlocus{No. 4, p. 2, leaf 83, left column. Ps 24(25):3--4.}

\begin{latinproper}
Univérsi, qui te exspéctant, non confundéntur, Dómine. V. Vias tuas, Dómine, notas fac mihi: et sémitas tuas édoce me.
\end{latinproper}

\begin{englishwitness}{Challoner, with historical liturgical forms}
None of them that wait on thee shall be confounded. V. Shew, O Lord, thy ways to me, and teach me thy paths.
\end{englishwitness}

\begin{englishwitness}{Shepherd, Gradual response with final vocative}
None of them that wait on thee shall be confounded, O Lord.
\end{englishwitness}

\begin{englishwitness}{Lasance and Walsh, hand-missal English}
None of them that wait on Thee shall be confounded. V. Show, O Lord, Thy ways to me, and teach me Thy paths.
\end{englishwitness}

\subsection{Alleluia}
\label{proper-alleluia}
\properlocus{No. 5, p. 2, leaf 83, left column. Ps 84:8 (85:7).}

\begin{latinproper}
Allelúia, allelúia. V. Osténde nobis, Dómine, misericórdiam tuam: et salutáre tuum da nobis. Allelúia.
\end{latinproper}

\begin{englishwitness}{Challoner, with historical liturgical forms}
Alleluia, alleluia. V. Shew us, O Lord, thy mercy; and grant us thy salvation. Alleluia.
\end{englishwitness}

\begin{englishwitness}{Lasance and Walsh, hand-missal English}
Alleluia, alleluia. V. Show us, O Lord, Thy mercy: and grant us Thy salvation. Alleluia.
\end{englishwitness}

\subsection{Gospel}
\label{proper-gospel}
\properlocus{No. 6, p. 2, leaf 83, both columns. Lk 21:25--33.}

\begin{latinproper}
Sequéntia sancti Evangélii secúndum Lucam. In illo témpore: Dixit Iesus discípulis suis: Erunt signa in sole, et luna, et stellis, et in terris pressúra géntium præ confusióne sónitus maris, et flúctuum: arescéntibus homínibus præ timóre et exspectatióne, quæ supervénient univérso orbi: nam virtútes cælórum movebúntur. Et tunc vidébunt Fílium hóminis veniéntem in nube cum potestáte magna, et maiestáte. His autem fíeri incipiéntibus, respícite, et leváte cápita vestra: quóniam appropínquat redémptio vestra. Et dixit illis similitúdinem: Vidéte ficúlneam, et omnes árbores: cum prodúcunt iam ex se fructum, scitis quóniam prope est æstas. Ita et vos cum vidéritis hæc fíeri, scitóte quóniam prope est regnum Dei. Amen dico vobis, quia non præteríbit generátio hæc, donec ómnia fiant. Cælum et terra transíbunt: verba autem mea non transíbunt.
\end{latinproper}

\begin{englishwitness}{Challoner, with historical liturgical forms}
Continuation of the holy Gospel according to St. Luke. At that time, Jesus said to His disciples: There shall be signs in the sun and in the moon and in the stars; and upon the earth distress of nations, by reason of the confusion of the roaring of the sea, and of the waves: Men withering away for fear and expectation of what shall come upon the whole world. For the powers of heaven shall be moved. And then they shall see the Son of man coming in a cloud, with great power and majesty. But when these things begin to come to pass, look up and lift up your heads, because your redemption is at hand. And he spoke to them a similitude. See the fig tree and all the trees: When they now shoot forth their fruit, you know that summer is nigh; So you also, when you shall see these things come to pass, know that the kingdom of God is at hand. Amen, I say to you, this generation shall not pass away till all things be fulfilled. Heaven and earth shall pass away: but my words shall not pass away.
\end{englishwitness}

\begin{englishwitness}{Lasance and Walsh, hand-missal English}
Continuation of the holy Gospel according to St. Luke. At that time, Jesus said to His disciples: There shall be signs in the sun, and in the moon, and in the stars: and upon the earth distress of nations, by reason of the confusion of the roaring of the sea and of the waves, men withering away for fear and expectation of what shall come upon the whole world. For the powers of heaven shall be moved; and then they shall see the Son of man coming in a cloud with great power and majesty. But when these things begin to come to pass, look up and lift up your heads, because your redemption is at hand. And He spoke to them a similitude: See the fig-tree, and all the trees: when they now shoot forth their fruit, you know that summer is nigh; so you also, when you shall see these things come to pass, know that the kingdom of God is at hand. Amen I say to you, this generation shall not pass away, till all things be fulfilled. Heaven and earth shall pass away, but My words shall not pass away.
\end{englishwitness}

\subsection{Offertory}
\label{proper-offertory}
\properlocus{No. 7, p. 2, leaf 83, right column. Ps 24(25):1--3.}

\begin{latinproper}
Ad te levávi ánimam meam: Deus meus, in te confído, non erubéscam: neque irrídeant me inimíci mei: étenim univérsi, qui te exspéctant, non confundéntur.
\end{latinproper}

\begin{englishwitness}{Challoner, with historical liturgical forms}
To thee, O Lord, have I lifted up my soul. In thee, O my God, I put my trust; let me not be ashamed. Neither let my enemies laugh at me: for none of them that wait on thee shall be confounded.
\end{englishwitness}

\begin{englishwitness}{Lasance and Walsh, hand-missal English}
To Thee have I lifted up my soul: in Thee, O my God, I put my trust; let me not be ashamed: neither let my enemies laugh at me: for none of them that wait on Thee shall be confounded.
\end{englishwitness}

\subsection{Secret}
\label{proper-secret}
\properlocus{No. 8, p. 2, leaf 83, right column.}

\begin{latinproper}
Hæc sacra nos, Dómine, poténti virtúte mundátos, ad suum fáciant purióres veníre princípium. Per Dóminum nostrum Iesum Christum, Fílium tuum: Qui tecum vivit et regnat in unitáte Spíritus Sancti, Deus: per ómnia sǽcula sæculórum. R. Amen.
\end{latinproper}

\begin{englishwitness}{Lasance and Walsh, Sunday text}
Cleansing us by their mighty power, may these Holy Mysteries, O Lord, make us come more pure before Thee who art their author. Through our Lord.
\end{englishwitness}

\begin{englishwitness}{Full conclusion, Prayer Conclusions, p. 819}
Through Our Lord Jesus Christ, Thy Son, Who liveth and reigneth with Thee in the unity of the Holy Ghost, God, world without end. R. Amen.
\end{englishwitness}

\subsection{Preface of the Most Holy Trinity}
\label{proper-preface}
\properlocus{Appointment after no. 8, p. 2. Thanksgiving, no. 1082, p. 293, leaf 374.}

The Secret ends with the first exchange printed here. At Mass its ending and Amen are heard once, followed by the greeting.

\begin{latinproper}
V. Per omnia saecula saeculorum. R. Amen.

V. Dominus vobiscum. R. Et cum spiritu tuo.

V. Sursum corda. R. Habemus ad Dominum.

V. Gratias agamus Domino Deo nostro. R. Dignum et iustum est.
\end{latinproper}

\begin{englishwitness}{Lasance and Walsh, dialogue}
P. World without end. R. Amen.

P. The Lord be with you. R. And with thy spirit.

P. Lift up your hearts. R. We have them lifted up unto the Lord.

P. Let us give thanks to the Lord our God. R. It is meet and just.
\end{englishwitness}

\begin{latinproper}
Vere dignum et iustum est, aequum et salutare, nos tibi semper et ubique
gratias agere: Domine, sancte Pater, omnipotens aeterne Deus: Qui cum
unigenito Filio tuo, et Spiritu Sancto, unus es Deus, unus es Dominus:
non in unius singularitate personae, sed in unius Trinitate substantiae.
Quod enim de tua gloria, revelante te, credimus, hoc de Filio tuo, hoc
de Spiritu Sancto, sine differentia discretionis sentimus. Ut in
confessione verae sempiternaeque Deitatis, et in personis proprietas,
et in essentia unitas, et in maiestate adoretur aequalitas. Quam laudant
Angeli atque Archangeli, Cherubim quoque ac Seraphim: qui non cessant
clamare cotidie, una voce dicentes:
\end{latinproper}

\begin{englishwitness}{Lasance and Walsh, thanksgiving}
It is truly meet and just, right and profitable, for us, at all times, and
in all places, to give thanks to Thee, O Lord, the holy One, the Father
almighty, the everlasting God: Who, together with Thine only-begotten Son
and the Holy Ghost, art one God, one Lord, not in the singleness of one
Person, but in the Trinity of one substance. For that which, according to
Thy revelation, we believe of Thy glory, the same we believe of Thy Son,
the same of the Holy Ghost, without difference or distinction; so that
in the confession of one true and eternal Godhead we adore distinctness
in persons, oneness in essence, and equality in majesty: Which the angels
praise, and the archangels, the cherubim also and the seraphim, who cease
not, day by day crying out with one voice, to repeat:
\end{englishwitness}

The common Sanctus follows.

\begin{latinproper}
Sanctus, Sanctus, Sanctus Dominus Deus Sabaoth. Pleni sunt caeli et terra
gloria tua. Hosanna in excelsis. Benedictus qui venit in nomine Domini.
Hosanna in excelsis.
\end{latinproper}

\begin{englishwitness}{Lasance and Walsh, Sanctus}
Holy, holy, holy, Lord God of hosts. The heavens and the earth are full
of Thy glory. Hosanna in the highest. Blessed is He Who cometh in the
name of the Lord. Hosanna in the highest.
\end{englishwitness}

\subsection{Communion}
\label{proper-communion}
\properlocus{No. 9, p. 2, leaf 83, right column. Ps 84:13 (85:12).}

\begin{latinproper}
Dóminus dabit benignitátem: et terra nostra dabit fructum suum.
\end{latinproper}

\begin{englishwitness}{Challoner, with historical liturgical forms}
The Lord will give goodness: and our earth shall yield her fruit.
\end{englishwitness}

\subsection{Postcommunion}
\label{proper-postcommunion}
\properlocus{No. 10, p. 2, leaf 83, right column.}

\begin{latinproper}
Orémus. Suscipiámus, Dómine, misericórdiam tuam in médio templi tui: ut reparatiónis nostræ ventúra solémnia cóngruis honóribus præcedámus. Per Dóminum nostrum.
\end{latinproper}

\begin{englishwitness}{Lasance and Walsh, Sunday text}
May we receive Thy mercy, O Lord, in the midst of Thy temple, that with due reverence we may prepare for the coming festival of our redemption. Through our Lord.
\end{englishwitness}

\begin{englishwitness}{Full conclusion, Prayer Conclusions, p. 819}
Through Our Lord Jesus Christ, Thy Son, Who liveth and reigneth with Thee in the unity of the Holy Ghost, God, world without end. R. Amen.
\end{englishwitness}
